\begin{abstract}
Extra inference structure can help a language model search, verify, and revise,
but those actions also consume the budget they are meant to use well. This
protocol preregisters a conditional test: verified search should be worse than a
one-call monolith at a low token budget and better at a high token budget. The
primary datasets are FinQA and TAT-QA, prepared from pinned local snapshots with
one selected question per source document, hidden labels, executable gold
derivations, and document-disjoint splits. The systems are a monolith, verified
search, and an exploratory unverified-search ablation, all using exactly
\texttt{gpt-5.4-mini}. Low, middle, and high tiers provide action-backed limits
on retrieval, planning, candidates, repair, and total authoritative prompt plus
completion tokens. Confirmation is an intersection-union test requiring both a
negative low-tier paired effect and a positive high-tier paired effect; equality
is never a crossover. FinanceComplexQA remains exploratory behind scorer,
lineage, leakage, oracle-evidence, retrieval, and orchestration boundaries. This
is a protocol-only artifact. No empirical conclusion is available, and the
hypothesis is currently neither supported nor cleanly falsified.
\end{abstract}
